 \documentclass{article}
\usepackage{tikz}
\usetikzlibrary{angles, quotes}
\usepackage{amssymb}
\usepackage{pgfplots}
\pgfplotsset{width=10cm,compat=1.9}
\usepgfplotslibrary{statistics}
\usepackage[left=2.54cm,right=2.54cm,top=2.54cm,bottom=2.54cm]{geometry}
\usepackage{tasks} %Paquete necesario para la producción de items horizontales para algunos ejercicios de tarea empleando el comando /begin{multicols}}{#}.../end{multicols} para ayudarnos a reducir las páginas
\usepackage{pdflscape} %comando que permite cambiar a orientación horizontal las hojas%
\usepackage{soul} %Paquete para permitir la compilación de acentos usando el comando {\hl{}}}\\
\sethlcolor{yellow(munsell)} %comando para definir el color para el subrayado usando el comando \hl
\usepackage[utf8]{inputenc}
\usepackage[spanish]{babel}
\usepackage{icomma}
\usepackage{ragged2e}
\usepackage{siunitx}
\usepackage{url}
\usepackage[colorlinks=true, urlcolor=blue,  linkcolor=black, citecolor=green]{hyperref}
\usepackage{pdfpages}
\usepackage{blindtext} %Paquete que genera texto ficticio (dummy text) con el comando: \blindtext
\usepackage{cancel} %in the preamble gives you four different modes of striking through: \cancel{text to cancel} draws a diagonal line (slash) through its argument, \bcancel{text to cancel} uses the negative slope (a backslash), \xcancel{text to cancel} draws an X (actually \cancel plus \bcancel), \cancelto{〈value〉}{〈expression〉} draws a diagonal arrow through the 〈expression〉pointing to the 〈value〉 (math-mode only)
\usepackage{csquotes}
\usepackage{afterpage}
\usepackage{parskip} 
\usepackage{float}
\usepackage{enumitem}
\usepackage{multicol}%Paquete que permite la creación aislada de columnas de texto. De esta forma se reduce la cantidad de páginas en nuestro documento
\usepackage{lipsum} %Paquete que genera texto de relleno al igual que \usepackage{blindtext}, se usa con el comando \lipsum[#-#] los #(hashtags) delimitan el número de páginas que deseamos ocupar del paquete lipsum para usarlos como relleno. 

% Definir el comando personalizado para la nota
\newcommand{\mynote}[1]{%
\begin{tikzpicture}[baseline=-0.75ex]
\node[draw, circle, fill=Apple Green!60, inner sep=2pt] (note) {\textbf{Nota:}};
\end{tikzpicture}%
\ \textit{#1}%
}

\newenvironment{Figura}
{\par\medskip\noindent\minipage{\linewidth}}
{\endminipage\par\medskip}
\usepackage{caption}
\usepackage[
backend=biber,
sorting=none,
url=true
]{biblatex} %bibliografía
\addbibresource{biblio.bib}
% Establecer el espacio entre las entradas de la bibliografía
\setlength{\bibitemsep}{1\baselineskip} % Puedes ajustar el valor según tus preferencias
\usepackage{amsmath,amsthm,amssymb,amsfonts}
\usepackage{pifont} %Permite la compilación del comando \xmark: es el símbolo de la tachita
\usepackage{mathtools} %permite la compilación de símbolos para matrices
\usepackage{empheq} %Paquete que se relaciona con \usepackage[most]{tcolorbox} para la creación de cuadros/cajas de colores para delimitar los resultados matemáticos o líneas de texto usando la línea de comando: \begin{empheq}[box={\mymath[colback=orange(webcolor)!70,drop lifted shadow]}]{equation*} \end{empheq}
\usepackage[most]{tcolorbox} %Paquete que se relaciona con \usepackage{empheq} para la creación de cuadros/cajas de colores para delimitar los resultados matemáticos o líneas de texto
\renewcommand{\qedsymbol}{$\blacksquare$}
\usetikzlibrary{positioning,decorations.pathreplacing} %Paquete que permite la compilación de llaves para cuadros sinópticos (brace diagram) 
\usepackage{schemata} %The schemata package is designed just for these brace diagrams
\usepackage{booktabs}

\newcommand\AB[2]{\schema{\schemabox{#1}}{\schemabox{#2}}} %comando o paquete necesario para crear cuadros sinópticos

\newtcbox{\mymath}[1][]{%
nobeforeafter, math upper, tcbox raise base,
enhanced, colframe=blue!30!black,
colback=blue!30, boxrule=1pt,
#1} %Comando importante para encasillar los resultados matemáticos en cajas de diferentes colores, se relaciona con los paquetes \usepackage{empheq} y \usepackage[most]{tcolorbox} 

\newenvironment{sysmatrix}[1]
{\left(\begin{array}{@{}#1@{}}}
{\end{array}\right)}
\newcommand{\ro}[1]{%
\xrightarrow{\mathmakebox[\rowidth]{#1}}%
}
\newlength{\rowidth}% row operation width
\AtBeginDocument{\setlength{\rowidth}{3em}}  %Comando importante para laproducción de lineas de operaciones entre matrices, método de Gauss o Gauss Jordan se relaciona con \newenvironment{sysmatrix}

%formato para cambiar el horario a español
\usepackage[spanish]{datetime2}
\DTMsetdatestyle{spanish}

\renewcommand{\today}{\DTMdisplaydate{\the\year}{\the\month}{\the\day }{-1}}
%%%%%%%%%%%%%%%%%%%%%%%%%%%%%%%%%%%%%%%%%

\usepackage{datetime}
\newcommand{\mycurrenttime}{\xxivtime}

%%%%%%%%%%%%%%%%%%%%%%%%%%%%%%%%%%%%%%%%%%%%%%%%%%%%%%%%%%%%%%%%%%%%%%%%%%%%%%%%%%%%%%%%%%%%%%%%%%%%%%%%%%%%%%%%%%%%%%%%%%%%%%%%%%%%%%%%%%%%%%%%%


%%%%%%%%%%%%%Caja de color para el título%%%%%%%%%%%%%%%%%%%%%%%%%%%%%%%%%%%%%%%%%%%%%%%%%%%%%%

\definecolor{myframecolor}{RGB}{1, 45, 75} % Prussian Blue
\definecolor{myboxcolor}{RGB}{0, 107, 92} % Apple Green

\newtcolorbox{mybox}{
enhanced,
colback=myboxcolor!7, % Color de fondo del cuadro
colframe=myframecolor, % Color del marco del cuadro
arc=0pt,
boxrule=1pt,
borderline west={2mm}{-10mm}{myframecolor}, % Borde en el lado izquierdo
sharp corners=southwest,
width=\linewidth,
before=\par\vspace{\bigskipamount}, % Espacio antes del cuadro
after=\par\vspace{\bigskipamount} % Espacio después del cuadro
}


\newcommand{\euler}{e} %Comando para producir letra e de euler
\newenvironment{remark}{\par\vfill\footnotesize % Vertical white space above the remark and smaller font size
\begin{list}{}{
	\leftmargin=80pt % Indentation on the left
	\rightmargin=60pt}\item\ignorespaces % Indentation on the right
\makebox[-2.5pt]{\begin{tikzpicture}[overlay]
		\node[draw=Horizon!60,line width=2.5pt,circle,fill=Horizon!25,font=\sffamily\bfseries,inner sep=4pt,outer sep=0pt] at (-19pt,5pt){\textcolor{Horizon}{Nota}};\end{tikzpicture}} % Blue Nota in a circle
\advance\baselineskip -1pt}{\end{list}\vskip5pt} % Tighter line spacing and white space after remark
\usepackage{graphicx}

\usepackage{titling}

\input{colores.tex}

%%Producir signos de cita textual``Asignatura''

\newcommand{\xmark}{\ding{55}} %%%comando que genera la tachita

\newenvironment{MyColorPar}[1]{%
\leavevmode\color{#1}\ignorespaces%
}{%
}%

%%%%%%%%%%%%%%%%%%% Título de la tarea, Nombre de alumno

\title{Clase \#30 - Electrodinámica}
\author{\emph{Julio Alfredo Ballinas García} $\boldsymbol{\mid}$ 202107583}
\date{\today}

%%%%%%%%%%%%%%%%%%% Datos de la Materia

\usepackage{fancyhdr}
\fancypagestyle{plain}{%  the preset of fancyhdr 
\fancyhf{} % clear all header and footer fields
\fancyfoot[R]{
\includegraphics[width=2cm]{LogoFCFMBUAP (1).png}}
\fancyfoot[L]{{\bfseries{\thedate{}}} a las {\bfseries{\mycurrenttime{}}} horas{} (GMT-6, H. Puebla de Zaragoza, Pue)}
\fancyhead[L]{Electrodinámica}
\fancyhead[R]{\theauthor}
}
\makeatletter
\def\@maketitle{%
\newpage
\null
\vskip 1em%
\begin{center}%
\let \footnote \thanks
{\LARGE \@title \par}%
\vskip 1em%
%{\large \@date}%
\end{center}%
\par
\vskip 1em}
\makeatother

\usepackage{lipsum}  
\usepackage{cmbright}


%%%%%%%%%%%%%%%%%%%%%%%%%%%%%%%%%%%%%%%%%%%%%%%%%%%%%%%%%%%%%%%%%%%%%%%%%%%%%%%%%%%%%%%%%%%%%%%%%%%%%%%%%%%%%%%%%%%%%%%%%%%%%%%%%%%%%%%%%%%%%%%%%%%%%%%%%%%%%%%%%%%%%%%%%%%%%%%%%%%%%%%%%%%%%%%%%%%%%%

\begin{document}



\maketitle


%%%%%%%%%%%%%%%%%%% Datos del profesor, campus y aula

\begin{mybox}
\noindent\begin{tabular}{@{}ll}
	{\bfseries{Profesor}} & Dr. Iván Fuentecilla Cárcamo\\
	\textcolor{prussianblue}{\bfseries{Campus}} / \textcolor{trueblue}{\bfseries{Facultad}}  & \textcolor{prussianblue}{\bfseries{C.U. BUAP}} /  \textcolor{trueblue}{\bfseries{FCFM}} \\
	\textcolor{prussianblue}{\bfseries{Edificio}} /  \textcolor{trueblue}{\bfseries{Salón}}    & 	\textcolor{prussianblue}{\bfseries{EMA4}} /  \textcolor{trueblue}{\bfseries{409}}   
\end{tabular} 
\end{mybox} \vspace{1cm}




	





\begin{MyColorPar}{orange(colorwheel)}

\end{MyColorPar}


\tableofcontents \newpage

\begin{center}
\begin{tikzpicture}
	% Dibuja círculos rellenos en diferentes colores
	\fill[bittersweet] (0,0) circle (0.17cm);
	\fill[Apple Green] (0.6,0) circle (0.17cm);
	\fill[trueblue] (1.2,0) circle (0.17cm);
	\fill[Aguamarina] (1.8,0) circle (0.17cm);
	\fill[Sun] (2.4,0) circle (0.17cm);
	\fill[orange(colorwheel)] (3.0,0) circle (0.17cm);
	\fill[Pesto] (3.6,0) circle (0.17cm);
	\fill[blue-violet] (4.2,0) circle (0.17cm);
	\fill[ticklemepink] (4.8,0) circle (0.17cm);
	\fill[Gallery] (5.4,0) circle (0.17cm);
	\fill[Mantle] (6.0,0) circle (0.17cm);
	\fill[Nevada] (6.6,0) circle (0.17cm);
	\fill[onyx] (7.2,0) circle (0.17cm);
	% Agrega más círculos o modifica los colores según sea necesario
\end{tikzpicture}
\end{center}  \vspace{2cm}




\section{Unit 2: Electric fields in dielectric media}

Placas paralelas infinitas, parte instercial dielectrico también infinito, esta rellenando esa región

 \begin{figure}[H]
    \centering
    \includegraphics[width=0.7\textwidth]{1.png} % ajusta el ancho según convenga
    \caption{Descripción clara y concisa de la figura.}
    \label{fig:mi_figura} % opcional: para usar con \ref{fig:mi_figura}
\end{figure}

¿Cuáles serían las superficies?

Tendríamos las superficies $S_{1}$ y $S_{2}$.

Calculamos la densidad de la carga ligada ($\sigma$)

¿Cuál es el signo de las densidades de carga superficial?

NO confundirnos con los signos de las placas. Asociado a las sigmas tendríamos también un campo electrico llamado de \underline{polarización}.

¿En qué dirección va a apuntar ese campo asociado a las distribuciones de $\sigma$?

Eliminamos las lineas de carga asociados a las placas. entonces ódemos ver la dependecia de las $\sigma$. Sería de \textbf{derecha} a \textbf{izquierda}. Además estamos suponiendo que es un vector constante, es decir tendríamos un + $\mathcal{P}$ o $ -\mathcal{P}$, siendo p constante. 

De izquierda a derecha va el campo eléctrico de las placas.

No tenemos una dependencia espacial en la $\mathcal{P}$. 

La suposición fuerte es:

\[ \nabla \cdot \mathcal{P}  \]

Tenemos que considerar una expansión multiplar, es decir no llegamos a algo tan diferente.

Tenemos:


\[ \rho (\vec{r}) =0 , \quad \sigma(\vec{r}) = \vec{\mathcal{P}} \cdot \hat{n} \]


\subsection{Depolarizing electric field}

\[ \vec{E}_{d}\]

\[ \vec{E}_{m} = \vec{E}_{ext} + \vec{E}_{d} \]

Tenemos:

\[  \vec{E}_{d} = -  \frac{\vec{ \mathcal{P}  }  }{{\epsilon_{0}}  }\]

\subsection{Campo molecular}

\[  \vec{E}_{m} =\underbrace{ \vec{E}_{ext} + \vec{E}_{d}}_{\vec{E}} + \vec{E}_{S} + \vec{E}^{\prime}  \]

No aislamos el término $\vec{E}_{S}$


Definimos: 

\[  \vec{E}_{m} = \vec{E} + \vec{E}_{s} + \underbrace{\vec{E}^{\prime}}_{\text{Para redes cristalinas $\vec{E}=0$}} \]

Nos queda por determinar:

\[\vec{E}_{S}\]

Este término está asociado a la superficie, es decir:

\[  \sigma_{\mathcal{P}} = \vec{\mathcal{P}} \cdot \hat{n}\]

 \begin{figure}[H]
    \centering
	    \includegraphics[width=0.7\textwidth]{2.png} % ajusta el ancho según convenga
    \caption{Descripción clara y concisa de la figura.}
    \label{fig:mi_figura} % opcional: para usar con \ref{fig:mi_figura}
\end{figure}


La polarización eléctrica fue calculada como un límite.

¿Cómo se calcula un potencial asociado a la carga puntual?

\textbf{Recordatorio:} ¿Cómo se calcula el campo eléctrico asociado a la $\sigma_{\mathcal{P}}$?

 \begin{figure}[H]
    \centering
	    \includegraphics[width=0.7\textwidth]{3.png} % ajusta el ancho según convenga
    \caption{Descripción clara y concisa de la figura.}
    \label{fig:mi_figura} % opcional: para usar con \ref{fig:mi_figura}
\end{figure}

\subsection{Depolarizing electric field (vía Gauss)}

\textbf{Geometría y suposiciones.} Placas paralelas infinitas con un dieléctrico infinito entre ellas.
Polarización uniforme \(\vec{\mathcal P}= \mathcal P\,\hat{x}\) (sin dependencia espacial).

\begin{equation}
\rho_b \equiv -\nabla\cdot\vec{\mathcal P}=0 \qquad (\text{polarización uniforme})
\end{equation}

\begin{equation}
\sigma_b \equiv \vec{\mathcal P}\cdot\hat{n}
\end{equation}

Tomamos las superficies del dieléctrico:
\(S_1\) (cara izquierda, normal \(\hat{n}_1=-\hat{x}\)) y
\(S_2\) (cara derecha, normal \(\hat{n}_2=+\hat{x}\)).

\begin{equation}
\sigma_b(S_1)=\vec{\mathcal P}\cdot(-\hat{x})=-\mathcal P,
\qquad
\sigma_b(S_2)=\vec{\mathcal P}\cdot(+\hat{x})=+\mathcal P.
\end{equation}

\textbf{Condición de salto (ley de Gauss en una "pillbox" que cruza la interfaz):}
para una caja gaussiana delgada atravesando la superficie con normal \(\hat{n}\),
\begin{equation}
\hat{n}\cdot\big(\vec{E}_{\text{fuera}}-\vec{E}_{\text{dentro}}\big) \;=\; \frac{\sigma_{\text{total}}}{\varepsilon_0}.
\end{equation}

Aplicada a las \emph{cargas ligadas} (considerando sólo el campo que ellas producen):

\begin{equation}
\hat{n}_1\cdot\big(\vec{E}_{L}-\vec{E}_{\text{in}}\big)=\frac{\sigma_b(S_1)}{\varepsilon_0}
\quad\Rightarrow\quad
(-\hat{x})\cdot\big(\vec{E}_{L}-\vec{E}_{\text{in}}\big)=\frac{-\mathcal P}{\varepsilon_0},
\end{equation}

\begin{equation}
\hat{n}_2\cdot\big(\vec{E}_{R}-\vec{E}_{\text{in}}\big)=\frac{\sigma_b(S_2)}{\varepsilon_0}
\quad\Rightarrow\quad
(+\hat{x})\cdot\big(\vec{E}_{R}-\vec{E}_{\text{in}}\big)=\frac{+\mathcal P}{\varepsilon_0},
\end{equation}

donde \(\vec{E}_{\text{in}}\) es el campo entre las caras, y \(\vec{E}_{L},\vec{E}_{R}\) son los campos
a la izquierda y a la derecha del dieléctrico (exterior).

Para dos láminas infinitas con densidades \(\pm \mathcal P\), por simetría y superposición
el campo fuera se anula:
\begin{equation}
\vec{E}_{L}=\vec{0},\qquad \vec{E}_{R}=\vec{0}.
\end{equation}

Sustituyendo:

\begin{equation}
(-\hat{x})\cdot( -\vec{E}_{\text{in}} )=\frac{-\mathcal P}{\varepsilon_0}
\quad\Rightarrow\quad
\vec{E}_{\text{in}}\cdot\hat{x}=-\frac{\mathcal P}{\varepsilon_0},
\end{equation}

\begin{equation}
(+\hat{x})\cdot( -\vec{E}_{\text{in}} )=\frac{+\mathcal P}{\varepsilon_0}
\quad\Rightarrow\quad
\vec{E}_{\text{in}}\cdot\hat{x}=-\frac{\mathcal P}{\varepsilon_0},
\end{equation}

resultados consistentes que fijan el \emph{campo uniforme} entre las superficies:

\begin{equation}
\boxed{\, \vec{E}_{d} \;=\; -\,\frac{\mathcal P}{\varepsilon_0}\,\hat{x} \;=\; -\,\frac{\vec{\mathcal P}}{\varepsilon_0} \,}
\end{equation}

El signo muestra que el campo \(\vec{E}_{d}\) apunta \emph{de la cara con \(\sigma_b=+\mathcal P\)}
hacia la \emph{cara con \(\sigma_b=-\mathcal P\)} (opuesto a \(\vec{\mathcal P}\)).

Por superposición con el campo debido a las cargas libres en las placas:

\begin{equation}
\boxed{\, \vec{E}_{m} \;=\; \vec{E}_{\text{ext}} \;+\; \vec{E}_{d} \,}
\end{equation}

y, en dieléctrico lineal isotrópico \(\vec{\mathcal P}=\varepsilon_0\chi_e\,\vec{E}_{m}\),

\begin{equation}
\vec{E}_{d}=-\chi_e\,\vec{E}_{m}, 
\qquad
\vec{E}_{\text{ext}}=(1+\chi_e)\,\vec{E}_{m}=\frac{\varepsilon}{\varepsilon_0}\,\vec{E}_{m}.
\end{equation}
	



\end{document}